% fig-frame.tex -- the residual frame and the encoding fork.
% No data dependency. Annotations from RESULTS/target_divergence.json, cellline scope:
%   full 34.59 / shift 111.42 / residual 118.25 ; ceilings 0.7420 / 0.7420 / 0.7432.
% Layout notes:
%   * two subtraction ROWS, not an L-shape (the L put the generic term on a band that collided
%     with its own caption).
%   * the caption for the generic sits BELOW both rows, not between them -- between them it ran
%     under the divider and into the right-hand column.
%   * row labels are centred on profiles spaced 1.65cm apart; at 1.45cm "generic(c)" and
%     "residual r(d,c)" touch.
\begin{figure}[htbp]
\centering
\begin{tikzpicture}[
  x=1cm, y=1cm, font=\small,
  bar/.style={fill=black!55, draw=none},
  faintbar/.style={fill=black!22, draw=none},
  gene/.style={draw=black!35, rounded corners=1pt, inner sep=1.8pt, font=\ttfamily\scriptsize},
  shared/.style={gene, fill=black!5, text=black!45},
  differs/.style={gene, fill=black!16, text=black, draw=black!70, thick},
  lbl/.style={font=\scriptsize, text=black!70},
  strong/.style={font=\scriptsize, text=black!90},
  op/.style={font=\normalsize, text=black!75},
]
\newcommand{\prof}[4]{%
  \foreach \h [count=\i from 0] in {#3} {
    \fill[#4] ($(#1,#2)+(\i*3.0pt,0)$) rectangle ++(2.3pt,\h pt);
  }%
}

% ================================================ LEFT: two subtractions
\node[lbl]    at (0.52,3.34) {treated};
\node[lbl]    at (2.17,3.34) {control};
\node[strong] at (3.90,3.34) {shift $\shift(d,c)$};
\prof{0.10}{2.62}{17,12,9,14,7,10,5,8}{bar}
\node[op] at (1.48,2.92) {$-$};
\prof{1.75}{2.62}{10,8,6,11,4,7,3,6}{bar}
\node[op] at (3.13,2.92) {$=$};
\prof{3.40}{2.62}{7,4,3,3,3,3,2,2}{bar}

\node[strong] at (0.52,1.94) {shift};
\node[lbl]    at (2.17,1.94) {generic$(c)$};
\node[strong] at (3.95,1.94) {residual $\resid(d,c)$};
\prof{0.10}{1.22}{7,4,3,3,3,3,2,2}{bar}
\node[op] at (1.48,1.52) {$-$};
\prof{1.75}{1.22}{6,4,3,2,3,2,2,2}{faintbar}
\node[op] at (3.13,1.52) {$=$};
\prof{3.40}{1.22}{1,0,0,1,0,1,0,0}{bar}

\node[lbl, anchor=north, align=center] at (2.30,0.82)
  {generic$(c)$ is the mean over the \emph{other}\\drugs in this cell line};

\draw[black!25] (5.05,0.10) -- (5.05,3.60);

% ================================================ RIGHT: the encoding fork
\node[strong, anchor=west] at (5.35,3.44) {the \emph{same} measurement, written two ways};

\node[lbl, anchor=west] at (5.35,2.96) {\textbf{(a)} genes ranked by \emph{expression}};
\node[anchor=west] at (5.35,2.52) {%
  \tikz[baseline]{\node[shared]{ACTB};}\;\tikz[baseline]{\node[shared]{GAPDH};}\;%
  \tikz[baseline]{\node[differs]{CDKN1A};}\;\tikz[baseline]{\node[shared]{TPT1};}\;$\cdots$};
\node[anchor=west] at (5.35,2.06) {%
  \tikz[baseline]{\node[shared]{ACTB};}\;\tikz[baseline]{\node[shared]{GAPDH};}\;%
  \tikz[baseline]{\node[differs]{HMOX1};}\;\tikz[baseline]{\node[shared]{TPT1};}\;$\cdots$};
\node[strong, anchor=west] at (5.35,1.58)
  {only \textbf{34.6 / 200} tokens differ \;(\textbf{17\%})};

\node[lbl, anchor=west] at (5.35,1.00) {\textbf{(b)} genes ranked by \emph{signed residual}};
\node[anchor=west] at (5.35,0.56) {%
  \tikz[baseline]{\node[differs]{CDKN1A};}\;\tikz[baseline]{\node[differs]{MDM2};}\;%
  \tikz[baseline]{\node[gene, fill=black!3]{[DOWN]};}\;\tikz[baseline]{\node[differs]{CCNB1};}\;%
  $\cdots$};
\node[strong, anchor=west] at (5.35,0.08)
  {\textbf{118.2 / 200} differ \;(\textbf{59\%})};

% ================================================ the invariant
\draw[decorate, decoration={brace, amplitude=4pt}, black!55] (0.10,-0.55) -- (11.20,-0.55);
\node[strong, anchor=north] at (5.65,-0.73)
  {and the replicate-retrieval ceiling is \textbf{unchanged}: $0.742 \rightarrow 0.743$};

\end{tikzpicture}
\caption[The residual frame and the encoding fork]{\textbf{The information was always present; the
tokenisation discarded it.} Left: the drug-specific residual is built by two subtractions. The
control removes the cell's own state, leaving the \emph{shift}; the mean over the other drugs in the
same cell line removes the \emph{generic} program that every compound triggers, leaving what makes
this drug different. Right: that same measurement, written two ways. Ranking genes by expression (a)
puts housekeeping genes first, so only $34.6$ of the top $200$ tokens differ between two drugs in the
same context and cross-entropy has almost nothing to learn drug identity from; ranking by signed
residual (b) raises that to $118.2$. The bracket is the point of the figure --- the
replicate-retrieval ceiling is essentially identical under both encodings, so no information has been
added. What changed is how much of it the token sequence exposes to the loss.}
\label{fig:frame}
\end{figure}
